\documentclass[10pt,a4paper]{amsart}
\usepackage[margin=19mm]{geometry}
\usepackage{amsmath,amssymb,mathtools}
\usepackage[scr=rsfs,cal=euler]{mathalfa}
\usepackage{xfrac}
\usepackage[unicode,hidelinks,bookmarks=false]{hyperref}
\newcommand{\ie}{i.e.\ }
\newcommand{\cf}{cf.\ }
\newcommand{\resp}{respectively\ }
\newcommand{\Subsection}{\subsection}
\providecommand{\nobreakdash}{\nobreak\hbox{-}}
\setlength{\emergencystretch}{2em}
\multlinegap0pt
\usepackage{bm}
\usepackage{bbm}
\usepackage{dsfont}
\usepackage{xspace}
\usepackage{cancel}
\usepackage{mathtools}
\usepackage{amsthm}
\makeatletter
\@ifundefined{thm}{\newtheorem{thm}{Thm}}{}
\@ifundefined{lem}{\newtheorem{lem}[thm]{Lemma}}{}
\makeatother
\newcommand{\R}{\mathbb{R}}
\title[Stability of Equilibria in a Biofilm Reactor Model with Wall]{Stability of Equilibria in a Biofilm Reactor Model with Wall Attachment and Thermodynamic Growth Inhibition}
\author{Katerina Nik, Christoph Walker}
\date{Source: arXiv:2607.08498v1 (CC BY 4.0)}
\begin{document}
\maketitle

\noindent\emph{Source and licence.} Stability of Equilibria in a Biofilm Reactor Model with Wall Attachment and Thermodynamic Growth Inhibition. Version arXiv:2607.08498v1 (CC BY 4.0), distributed under the Creative Commons Attribution 4.0 International licence, as recorded in the arXiv licence metadata for this exact version and shown on the abstract page https://arxiv.org/abs/2607.08498v1. Full rights notice: licenses/arxiv-2607.08498-CC-BY-4.0.txt.\par\smallskip

\noindent\textit{Editorial scope.} Counted unit: the Lem whose source label is \texttt{P30} in the article (source0.tex lines 1058--1060, source label \texttt{P30}), delivered together with the article's own complete original proof of it (source0.tex lines 1062--1121, one \texttt{proof} environment of 60 lines). No mathematics is written, repaired or re-derived: statement and proof are the article's own text line for line. Required context retained uncounted: GAshooting (lines 290--310); ZZ (lines 675--683); Z1x (lines 676--678); shooting (lines 776--781); L21 (lines 844--853); 381 (lines 878--880); 382 (lines 882--884); poss (lines 886--888); L22 (lines 896--927); 400 (lines 970--975); L40 (lines 979--984); Bq (lines 1021--1023); L30 (lines 1033--1039). Every definition, display and label of those lines is retained unchanged. Omitted from this package and not redistributed: 1--289, 311--674, 679--775, 782--843, 854--877, 881--881, 885--885, 889--895, 928--969, 976--978, 985--1020, 1024--1032, 1040--1057, 1061--1061, 1122--1925. Nothing inside a retained excerpt is omitted or altered: every retained body is delivered exactly as the article's lines are written, so no omission receipt of any kind accompanies this package. Citations: the retained excerpts carry no citation command at all, so no bibliography entry of the article is transcribed and no reference is redistributed. Reference adaptation: none was needed and none was made; every reference inside the delivered unit resolves inside it, so no unresolved reference or citation is produced. Printed slips kept literal: no printed word, symbol, hypothesis or number of the counted statement or of its proof was altered. Layout and class: the article's own class is \texttt{amsart} with packages including times, amsmath, cancel, stmaryrd, graphicx, amsfonts, enumitem, amssymb, color, mathrsfs, hyperref, cleveref, enumerate, lmodern; the wrapper is the lane's pinned \texttt{amsart} editorial wrapper at 10pt on A4, which restates the theorem-like environments the retained text needs and adds the article's own macro definitions that the retained text needs (\texttt{\textbackslash R}), with extra wrapper packages (bm, bbm, dsfont, xspace, cancel, mathtools). The delivered numbering is the unit's own. Assets: no retained span contains a figure environment, a graphics-inclusion command, a PDF-page inclusion, a table or a TikZ picture; no image file is copied into this package, the package declares no asset dependency and no image file is redistributed with it. Licence: version arXiv:2607.08498v1 of this article is distributed under the Creative Commons Attribution 4.0 International licence, as recorded in the arXiv licence metadata for this exact version and shown on the abstract page https://arxiv.org/abs/2607.08498v1. A full rights notice is delivered at licenses/arxiv-2607.08498-CC-BY-4.0.txt. Characters: every retained excerpt is pure ASCII, so the delivered engine, XeLaTeX with its default Latin Modern text font, reports no missing character.
\begin{subequations}\label{GAshooting}
\begin{align}
&r\in  {\rm C}^1( \R_+^2,\R_+)\,,\qquad d\in {\rm C}^1(\R_+,\R_+)\label{r1}
\end{align}
satisfy
\begin{align}
r(0,y)=0\,,\quad \partial_2 r(x,y)< 0<\partial_1 r(x,y)\,,\quad (x,y)\in (0,\infty)^2\,,\label{r2}
\end{align}
where $\R_+=[0,\infty)$, and
\begin{align}
d'(x)\ge 0\,,\quad x\ge 0\,, \qquad  d(x)>0\,,\quad x>0\,. \label{d}
\end{align}
For the diffusion coefficients $\kappa_1$ and $\kappa_2$ of propionate and acetate, respectively, we assume that
\begin{align}\label{kappa}
\kappa_2>\kappa_1\,.
\end{align}
Finally, we always consider nonnegative initial values
\begin{equation}
\big(h^0,S_1^0,S_2^0,Q^0\big)\in\R_+^4\,.
\end{equation}
\end{subequations}

\begin{subequations}\label{ZZ}
\begin{align}\label{Z1x}
&k_2\ge k_1 r(S^*,0) 
\end{align}
and the uniform instability condition
\begin{align}\label{h2eff}
\Big(k_1\,r\big(S^*,\kappa_2^{-1}(\kappa_2-\kappa_1)S^*\big)-b\Big)\big(\alpha+k_2-k_1 r(S^*,0)\big)>\big(k_2-k_1 r(S^*,0)\big)d(0)\,.
\end{align}
\end{subequations}

\begin{subequations}\label{shooting}
\begin{align}
&\kappa_1\partial_z^2 {\rm c}(z)=R({\rm c}(z),h_0)\,, \quad z>0\,, \label{A1ss}\\
&{\rm c}(0)=\mu\,,\quad \partial_z {\rm c}(0)=0\,. \label{A4ss}
\end{align}
\end{subequations}

\begin{lem}\label{L21}
Assume~\eqref{GAshooting} and \eqref{Z1x}.
Let ${\rm c}$  be the solution to the shooting problem~\eqref{shooting}.
Then, there exists a unique function $h\in {\rm C}^1([0,S^*]\times [0,h_*])$   such that 
\begin{align}\label{id}
{\rm c}\big(h(\mu,h_0),\mu,h_0\big)=S\big(h(\mu,h_0)\big)\,,\quad (\mu,h_0)\in [0,S^*]\times [0,h_*]\,.
\end{align}
For fixed $h_0\in [0,h_*]$, the mapping $\mu\mapsto h(\mu,h_0)$ is strictly decreasing on $[0,S^*]$ 
with $h(0,h_0)=h_*$ and $h(S^*,h_0)=0$, while for fixed $\mu\in (0,S^*)$, the mapping $h_0\mapsto h(\mu,h_0)$ is strictly increasing on~$[0,h_*]$.
\end{lem}

\begin{align}\label{381}
\big[\partial_z{\rm c}(h(\mu,h_0),\mu,h_0)-S'(h(\mu,h_0)) \big]\partial_{\mu} h(\mu,h_0)=-\partial_\mu {\rm c}(h(\mu,h_0),\mu,h_0)\le -1
\end{align}

\begin{align}\label{382}
\big[\partial_z {\rm c}(h(\mu,h_0),\mu,h_0)-S'(h(\mu,h_0)) \big]\partial_{h_0} h(\mu,h_0)=-\partial_{h_0} {\rm c}(h(\mu,h_0),\mu,h_0)>0\,.
\end{align}

\begin{equation}\label{poss}
\partial_z{\rm c}(h(\mu,h_0),\mu,h_0)-S'(h(\mu,h_0))>0
\end{equation}

\begin{lem}\label{L22}
Assume~\eqref{GAshooting} and~\eqref{ZZ}.  Given $(z,\mu,h_0)\in [0,h_*]\times [0,S^*]\times [0,h_*]$ and $h\in [0,h_*]$ set
\[
w(z,\mu,h_0):=\partial_\mu{\rm c}(z,\mu,h_0)\,,\qquad
\psi(h):=\frac{k_2-k_1{\rm r}(S(h))}{\Delta(S(h))}\,,
\]
and
\begin{equation}\label{Mdef}
\begin{aligned}
M(z,\mu,h_0):=&-\big(k_1 R({\rm c}(z,\mu,h_0),h_0)-b\big)w(z,\mu,h_0)\\
&+k_1\kappa_1\big(\partial_z{\rm c}(z,\mu,h_0)-S'(h(\mu,h_0))\big)\partial_z w(z,\mu,h_0)\\
&+\psi(h(\mu,h_0))d(h(\mu,h_0))\,.
\end{aligned}
\end{equation}
Then, for fixed $(\mu,h_0)\in [0,S^*]\times [0,h_*]$, the function
$$
M(\cdot,\mu,h_0):[0,h_*]\to \R
$$ 
is strictly increasing. Moreover, the function
$M(0,\cdot,h_0):[0,S^*]\to\R$ with
\begin{equation}\label{M0}
M(0,\mu,h_0)=- k_1 R(\mu,h_0)+b +\frac{k_2-k_1{\rm r}\big(S(h(\mu,h_0))\big)}{\Delta\big(S(h(\mu,h_0))\big)}d\big(h(\mu,h_0)\big)\,,\quad \mu\in [0,S^*]\,,
\end{equation}
is strictly decreasing for each $h_0\in [0,h_*]$, and there exists (a unique) $\underline{\mu}(h_0)\in (0,S^*)$ such that 
$$
M(0,\underline{\mu}(h_0),h_0)=0\,.
$$ 
In particular,
\begin{equation}\label{M}
M(z,\mu,h_0)> M(0,\mu,h_0)\ge 0\,,\qquad z\in (0,h_*]\,,\quad \mu\in [0,\underline{\mu}(h_0)]\,,\quad h_0\in [0,h_*] \,.
\end{equation}
\end{lem}

\begin{equation}\label{400}
\begin{split}
B(\mu,h_0):=\, & k_1\kappa_1 \partial_z {\rm c}(h(\mu,h_0),\mu,h_0)-bh(\mu,h_0)\\
&\ -\frac{k_2-k_1 {\rm r}(S(h(\mu,h_0)))}{\Delta(S(h(\mu,h_0)))}d\big(h(\mu,h_0)\big)h(\mu,h_0)
\end{split}
\end{equation}

\begin{lem}\label{L40}
Assume~\eqref{GAshooting} and~\eqref{ZZ}. Let  $\underline{\mu}\in {\rm C}^1\big([0,h_*], (0,S^*)\big)$ be as in Lemma~\ref{L22}. Then, there exists a function $\mu_*\in {\rm C}^1([0,h_*])$ such that, for every
$h_0\in[0,h_*]$, one has $\mu_*(h_0)\in(0,\underline{\mu}(h_0)]$, and
$\mu=\mu_*(h_0)$ and $\mu=S^*$ are the only zeros of
$B(\cdot,h_0):[0,S^*]\to\R$.
\end{lem}

\begin{equation}\label{Bq}
\partial_\mu B(\mu,h_0)>0\,,\quad \mu\in [0,\underline{\mu}(h_0)]\,.
\end{equation}  

\begin{lem}\label{L30}
Assume~\eqref{GAshooting} and \eqref{Z1x}.  Define
\begin{equation*}
K(z,\mu,h_0):=k_1\kappa_1\big[\partial_z\partial_{h_0}{\rm c}(z,\mu,h_0)\partial_\mu{\rm c}(z,\mu,h_0)-\partial_z\partial_\mu{\rm c}(z,\mu,h_0)\partial_{h_0}{\rm c}(z,\mu,h_0)\big]
\end{equation*}
for $(z,\mu,h_0)\in[0,\infty)\times[0,S^*]\times[0,h_*]$. Then $K(z,\mu,h_0)\le 0$.
\end{lem}

\begin{lem}\label{P30}
Assume~\eqref{GAshooting} and~\eqref{ZZ}. Let $h=h(\mu,h_0)$ and $\mu_*=\mu_*(h_0)$ be the functions from Lemma~\ref{L21} and Lemma~\ref{L40}, respectively. Then, there exists a unique $h_\star\in (0,h_*)$ such that $h\big(\mu_*(h_\star),h_\star)=h_\star$.
\end{lem}

\begin{proof}
Set
$$
D(h_0):=h\big(\mu_*(h_0),h_0)-h_0\,,\quad h_0\in [0,h_*]\,,
$$
and note from Lemma~\ref{L21} that
$$
D(0)=h(\mu_*(0),0)>0\,,\qquad  D(h_*)=h(\mu_*(h_*),h_*)-h_*<0\,.
$$
Therefore, the assertion follows, provided that we can show that $D'(h_0)<0$ for $h_0\in[0,h_*]$, where
\begin{equation}\label{D}
D'(h_0)=\partial_\mu h(\mu_*(h_0),h_0)\mu_*'(h_0)+\partial_{h_0}h(\mu_*(h_0),h_0)-1\,.
\end{equation}
Differentiating $B(\mu_*(h_0),h_0)=0$ with respect to $h_0$ yields
\begin{equation}\label{D1}
\partial_\mu B(\mu_*(h_0),h_0)\mu_*'(h_0)=-\partial_{h_0}B(\mu_*(h_0),h_0)\,.
\end{equation}
Writing $B$ from~\eqref{400} in the form 
$$
B(\mu,h_0)=k_1\kappa_1\partial_z{\rm c}(h(\mu,h_0),\mu,h_0)-\phi(h(\mu,h_0)) 
$$
with
$$
\phi(h):=bh+\frac{k_2-k_1{\rm r}(S(h))}{\Delta(S(h))}d(h)h\,,
$$
it follows that
\begin{equation}\label{D2}
\begin{split}
\partial_\mu B(\mu,h_0)&=k_1 R\big({\rm c}(h(\mu,h_0),\mu,h_0),h_0\big) \partial_\mu h(\mu,h_0)\\
&\quad +
k_1\kappa_1 \partial_\mu\partial_z {\rm c}(h(\mu,h_0),\mu,h_0) - \phi'\big(h(\mu,h_0)\big)\partial_\mu h(\mu,h_0)
\end{split}
\end{equation}
and
\begin{equation}\label{D3}
\begin{split}
\partial_{h_0} B(\mu,h_0)&=k_1 R\big({\rm c}(h(\mu,h_0),\mu,h_0),h_0\big) \partial_{h_0} h(\mu,h_0)\\
&\quad +
k_1\kappa_1 \partial_{h_0}\partial_z {\rm c}(h(\mu,h_0),\mu,h_0) - \phi'\big(h(\mu,h_0)\big)\partial_{h_0} h(\mu,h_0)\,.
\end{split}
\end{equation}
Now, we infer from~\eqref{D} and~\eqref{D1} that
\begin{align*}
\partial_{\mu}B(\mu_*(h_0),h_0)\big(D'(h_0)+1\big)&=-\partial_{h_0}B(\mu_*(h_0),h_0) \partial_{\mu}h\big(\mu_*(h_0),h_0)\\
&\quad +\partial_{\mu}B(\mu_*(h_0),h_0) \partial_{h_0} h(\mu_*(h_0),h_0) 
\end{align*}
and thus, taking into account~\eqref{D2} and~\eqref{D3}, we obtain
\begin{align*}
\partial_{\mu}B(\mu_*(h_0),h_0)\big(D'(h_0)+1\big)&=-k_1\kappa_1 \partial_{h_0}\partial_z {\rm c}\big(h(\mu_*(h_0),h_0),\mu_*(h_0),h_0\big) \partial_{\mu} h\big(\mu_*(h_0),h_0\big)\\
&\quad +k_1\kappa_1 \partial_{\mu}\partial_z {\rm c}\big(h(\mu_*(h_0),h_0),\mu_*(h_0),h_0\big) \partial_{h_0} h\big(\mu_*(h_0),h_0\big)\,.
\end{align*}
Therefore, invoking~\eqref{381}-\eqref{382} we derive
\begin{align*}
\big[\partial_z{\rm c}&(h(\mu_*(h_0),h_0),\mu_*(h_0),h_0)-S'(h(\mu_*(h_0),h_0)) \big]\partial_{\mu}B(\mu_*(h_0),h_0)\big(D'(h_0)+1\big)\\
&=k_1\kappa_1 \partial_{h_0}\partial_z {\rm c}\big(h(\mu_*(h_0),h_0),\mu_*(h_0),h_0\big) \partial_{\mu} {\rm c}\big(h(\mu_*(h_0),h_0),\mu_*(h_0),h_0\big)\\
&\quad -k_1\kappa_1 \partial_{\mu}\partial_z {\rm c}\big(h(\mu_*(h_0),h_0),\mu_*(h_0),h_0\big) \partial_{h_0} {\rm c}\big(h(\mu_*(h_0),h_0),\mu_*(h_0),h_0\big)\\
&=K\big(h(\mu_*(h_0),h_0),\mu_*(h_0),h_0\big)\,.
\end{align*}
Consequently, since $K$ is negative according to Lemma~\ref{L30}, we conclude from~\eqref{poss} and~\eqref{Bq} that $D'(h_0)\le -1$ for $h_0\in [0,h_*]$.
\end{proof}

\end{document}
